\documentclass[10pt,a4paper]{amsart}
\usepackage{amsmath,amssymb,amsthm,graphicx}
\usepackage{verbatim}
\newenvironment{reference}{\comment}{\endcomment}
%\newenvironment{reference}{}{}
\newenvironment{slogan}{\comment}{\endcomment}
\newenvironment{history}{\comment}{\endcomment}

% For commutative diagrams we use Xy-pic
\usepackage[all]{xy}

% We use 2cell for 2-commutative diagrams.
\xyoption{2cell}
% Theorem environments.
%
\theoremstyle{plain}
\newtheorem{theorem}[subsection]{Theorem}
\newtheorem{proposition}[subsection]{Proposition}
\newtheorem{lemma}[subsection]{Lemma}

\theoremstyle{definition}
\newtheorem{definition}[subsection]{Definition}
\newtheorem{example}[subsection]{Example}
\newtheorem{exercise}[subsection]{Exercise}
\newtheorem{situation}[subsection]{Situation}

\theoremstyle{remark}
\newtheorem{remark}[subsection]{Remark}
\newtheorem{remarks}[subsection]{Remarks}

\numberwithin{equation}{subsection}

% Macros
%
\def\lim{\mathop{\mathrm{lim}}\nolimits}
\def\colim{\mathop{\mathrm{colim}}\nolimits}
\def\Spec{\mathop{\mathrm{Spec}}}
\def\Hom{\mathop{\mathrm{Hom}}\nolimits}
\def\Ext{\mathop{\mathrm{Ext}}\nolimits}
\def\SheafHom{\mathop{\mathcal{H}\!\mathit{om}}\nolimits}
\def\SheafExt{\mathop{\mathcal{E}\!\mathit{xt}}\nolimits}
\def\Sch{\mathit{Sch}}
\def\Mor{\mathop{\mathrm{Mor}}\nolimits}
\def\Ob{\mathop{\mathrm{Ob}}\nolimits}
\def\Sh{\mathop{\mathit{Sh}}\nolimits}
\def\NL{\mathop{N\!L}\nolimits}
\def\CH{\mathop{\mathrm{CH}}\nolimits}
\def\proetale{{pro\text{-}\acute{e}tale}}
\def\etale{{\acute{e}tale}}
\def\QCoh{\mathit{QCoh}}
\def\Ker{\mathop{\mathrm{Ker}}}
\def\Im{\mathop{\mathrm{Im}}}
\def\Coker{\mathop{\mathrm{Coker}}}
\def\Coim{\mathop{\mathrm{Coim}}}

% Boxtimes
%
\DeclareMathSymbol{\boxtimes}{\mathbin}{AMSa}{"02}

%
% Macros for moduli stacks/spaces
%
\def\QCohstack{\mathcal{QC}\!\mathit{oh}}
\def\Cohstack{\mathcal{C}\!\mathit{oh}}
\def\Spacesstack{\mathcal{S}\!\mathit{paces}}
\def\Quotfunctor{\mathrm{Quot}}
\def\Hilbfunctor{\mathrm{Hilb}}
\def\Curvesstack{\mathcal{C}\!\mathit{urves}}
\def\Polarizedstack{\mathcal{P}\!\mathit{olarized}}
\def\Complexesstack{\mathcal{C}\!\mathit{omplexes}}
% \Pic is the operator that assigns to X its picard group, usage \Pic(X)
% \Picardstack_{X/B} denotes the Picard stack of X over B
% \Picardfunctor_{X/B} denotes the Picard functor of X over B
\def\Pic{\mathop{\mathrm{Pic}}\nolimits}
\def\Picardstack{\mathcal{P}\!\mathit{ic}}
\def\Picardfunctor{\mathrm{Pic}}
\def\Deformationcategory{\mathcal{D}\!\mathit{ef}}
\usepackage{cleveref}
\crefname{originalreference}{original reference}{original references}
\Crefname{originalreference}{Original reference}{Original references}
\makeatletter
\newlabel{sites-modules-section-internal-hom}{{04TT}{source}{}{sourcepointer-1}{}}
\newlabel{sites-modules-section-internal-hom@cref}{{[originalreference][0][]04TT}{[1][0][]source}{}{}{}}
\makeatother
\crefname{originalreference}{original reference}{original references}
\Crefname{originalreference}{Original reference}{Original references}
\makeatletter
\newlabel{sites-modules-lemma-internal-hom}{{03EM}{source}{}{sourcepointer-2}{}}
\newlabel{sites-modules-lemma-internal-hom@cref}{{[originalreference][0][]03EM}{[1][0][]source}{}{}{}}
\makeatother
\crefname{originalreference}{original reference}{original references}
\Crefname{originalreference}{Original reference}{Original references}
\makeatletter
\newlabel{stacks-sheaves-lemma-lqc-colimits}{{06WM}{source}{}{sourcepointer-3}{}}
\newlabel{stacks-sheaves-lemma-lqc-colimits@cref}{{[originalreference][0][]06WM}{[1][0][]source}{}{}{}}
\makeatother
\crefname{originallemma}{original lemma}{original lemmas}
\Crefname{originallemma}{Original lemma}{Original lemmas}
\makeatletter
\newlabel{lemma-check-flat-comparison-on-etale-covering}{{0764}{source}{}{sourcepointer-4}{}}
\newlabel{lemma-check-flat-comparison-on-etale-covering@cref}{{[originallemma][0][]0764}{[1][0][]source}{}{}{}}
\makeatother
\crefname{originallemma}{original lemma}{original lemmas}
\Crefname{originallemma}{Original lemma}{Original lemmas}
\makeatletter
\newlabel{lemma-adjoint}{{0778}{source}{}{sourcepointer-5}{}}
\newlabel{lemma-adjoint@cref}{{[originallemma][0][]0778}{[1][0][]source}{}{}{}}
\makeatother
\crefname{originalreference}{original reference}{original references}
\Crefname{originalreference}{Original reference}{Original references}
\makeatletter
\newlabel{sites-modules-lemma-internal-hom-adjoint-tensor}{{03EO}{source}{}{sourcepointer-6}{}}
\newlabel{sites-modules-lemma-internal-hom-adjoint-tensor@cref}{{[originalreference][0][]03EO}{[1][0][]source}{}{}{}}
\makeatother
\title{Internal hom from a finitely presented module into a qc module is lqc-fbc and has the tensor adjoint}
\author{The Stacks Project Authors}
\date{stacks-project:stacks-cohomology.tex:local-snapshot-20261012; GNU FDL 1.2 or later}
\begin{document}
\maketitle
\noindent\textit{Editorial note.} All delivered mathematical text is the author's own, copied verbatim from the pinned Stacks Project chapter file stacks-cohomology.tex; no statement or proof step was rewritten, reordered or removed.\par\medskip
\noindent\textit{Editorial note.} The counted span cites the lemmas recorded in assembly.external\_references with canonical tags as navigation pointers.\par\medskip
\noindent\textit{Allowed standard prerequisites (editorial).}\begin{itemize}
\item Internal hom sheafifies. Modules on Sites, Lemma sites-modules-lemma-internal-hom. fppf sheaf property.
\item fbc of Hom from fp modules. Cohomology of Algebraic Stacks, Lemma stacks-cohomology-lemma-check-flat-comparison-on-etale-covering (6). Flat base change property.
\item Internal hom tensor adjunction. Modules on Sites, Lemma sites-modules-lemma-internal-hom-adjoint-tensor. Final universal property.
\end{itemize}\medskip
\ifdefined\hypertarget\hypertarget{sourcepointer-1}{}\fi
\noindent\textit{Source reference.} Original reference 04TT: sources/stacks-chow/tags.txt, source lines 6266--6266. Cross-chapter standard prerequisite cited by the author inside the counted span; it is used as a pointer and its own proof is not part of this item.\par\medskip
\ifdefined\hypertarget\hypertarget{sourcepointer-2}{}\fi
\noindent\textit{Source reference.} Original reference 03EM: sources/stacks-chow/tags.txt, source lines 4461--4461. Cross-chapter standard prerequisite cited by the author inside the counted span; it is used as a pointer and its own proof is not part of this item.\par\medskip
\ifdefined\hypertarget\hypertarget{sourcepointer-3}{}\fi
\noindent\textit{Source reference.} Original reference 06WM: sources/stacks-chow/tags.txt, source lines 8821--8821. Cross-chapter standard prerequisite cited by the author inside the counted span; it is used as a pointer and its own proof is not part of this item.\par\medskip
\ifdefined\hypertarget\hypertarget{sourcepointer-4}{}\fi
\noindent\textit{Source reference.} Original lemma 0764: sources/stacks-chow/tags.txt, source lines 9154--9154. Cross-chapter standard prerequisite cited by the author inside the counted span; it is used as a pointer and its own proof is not part of this item.\par\medskip
\ifdefined\hypertarget\hypertarget{sourcepointer-5}{}\fi
\noindent\textit{Source reference.} Original lemma 0778: sources/stacks-chow/tags.txt, source lines 9193--9193. Cross-chapter standard prerequisite cited by the author inside the counted span; it is used as a pointer and its own proof is not part of this item.\par\medskip
\ifdefined\hypertarget\hypertarget{sourcepointer-6}{}\fi
\noindent\textit{Source reference.} Original reference 03EO: sources/stacks-chow/tags.txt, source lines 4463--4463. Cross-chapter standard prerequisite cited by the author inside the counted span; it is used as a pointer and its own proof is not part of this item.\par\medskip
\begin{lemma}
\label{lemma-internal-hom-fp-into-qcoh}
Let $\mathcal{X}$ be an algebraic stack. Let $\mathcal{F}$ be an
$\mathcal{O}_\mathcal{X}$-module of finite presentation and let
$\mathcal{G}$ be a quasi-coherent $\mathcal{O}_\mathcal{X}$-module.
The internal homs
$\SheafHom_{\mathcal{O}_\mathcal{X}}(\mathcal{F}, \mathcal{G})$
computed in
$\textit{Mod}(\mathcal{X}_\etale, \mathcal{O}_\mathcal{X})$ or
$\textit{Mod}(\mathcal{O}_\mathcal{X})$ agree and the common value
is an object of $\textit{LQCoh}^{fbc}(\mathcal{O}_\mathcal{X})$.
The quasi-coherent module
$
hom(\mathcal{F}, \mathcal{G}) =
Q(\SheafHom_{\mathcal{O}_\mathcal{X}}(\mathcal{F}, \mathcal{G}))
$
has the following universal property
$$
\Hom_\mathcal{X}(\mathcal{H}, hom(\mathcal{F}, \mathcal{G})) =
\Hom_\mathcal{X}(\mathcal{H} \otimes_{\mathcal{O}_\mathcal{X}} \mathcal{F},
\mathcal{G})
$$
for $\mathcal{H}$ in $\QCoh(\mathcal{O}_\mathcal{X})$.
\end{lemma}

\begin{proof}
The construction of
$\SheafHom_{\mathcal{O}_\mathcal{X}}(\mathcal{F}, \mathcal{G})$ in
Modules on Sites, Section \ref{sites-modules-section-internal-hom}
depends only on $\mathcal{F}$ and $\mathcal{G}$ as presheaves
of modules; the output $\SheafHom$ is a sheaf for the fppf
topology because $\mathcal{F}$ and $\mathcal{G}$ are assumed
sheaves in the fppf topology, see
Modules on Sites, Lemma \ref{sites-modules-lemma-internal-hom}.
By Sheaves on Stacks, Lemma \ref{stacks-sheaves-lemma-lqc-colimits}
we see that $\SheafHom_{\mathcal{O}_\mathcal{X}}(\mathcal{F}, \mathcal{G})$
is locally quasi-coherent.
By Lemma \ref{lemma-check-flat-comparison-on-etale-covering}
we see that $\SheafHom_{\mathcal{O}_\mathcal{X}}(\mathcal{F}, \mathcal{G})$
has the flat base change property. Hence
$\SheafHom_{\mathcal{O}_\mathcal{X}}(\mathcal{F}, \mathcal{G})$
is an object of $\textit{LQCoh}^{fbc}(\mathcal{O}_\mathcal{X})$
and it makes sense to apply the functor $Q$ of Lemma \ref{lemma-adjoint}
to it. By the universal property of $Q$ we have
$$
\Hom_\mathcal{X}(\mathcal{H},
Q(\SheafHom_{\mathcal{O}_\mathcal{X}}(\mathcal{F}, \mathcal{G}))) =
\Hom_\mathcal{X}(\mathcal{H},
\SheafHom_{\mathcal{O}_\mathcal{X}}(\mathcal{F}, \mathcal{G}))
$$
for $\mathcal{H}$ quasi-coherent, hence the displayed formula 
of the lemma follows from Modules on Sites, Lemma
\ref{sites-modules-lemma-internal-hom-adjoint-tensor}.
\end{proof}

\end{document}
